\documentclass{article}
\usepackage[utf8]{inputenc}
\usepackage[portuguese]{babel}

\title{Cálculo Numérico - Notas de Aula}
\author{Giulia Leano}
\date{10 de setembro de 2026}

\begin{document}

\maketitle

\section{Bibliotecas}
O professor comentou no início da aula que não iria passar por todas os capítulos do pdf, mas sim pelos mais importantes.
O primeiro capítulo que vimos, no caso, que de fato iremos usar, são as bibliotecas. As bibliotecas são 
coleções de módulos e funções reutilizáveis que simplificam o desenvolvimento de um programa, sem a necessidade de reescrever tudo do zero. A primeira biblioteca vista foi a math, onde contém funções matemáticas. 
No caso dela já estiver baixada, podemos chamá-la para a IDE com o comando import math e se ele não estiver baixada, podemos instalar no terminal utilizando o comando pip install math
\end{document}.